
En
effet, (ii) implique qu'il existe un diviseur effectif $D$ tel que $N^{\otimes n}\simeq{O}_X(D)$, et $\langle c_1(N)^2\rangle>0$ implique que $D>0$ ; d'o\`u (i).

Pour voir que (i') implique (ii), on peut \'evidemment supposer que $H={O}_X(1)$ est tr\`es ample ; alors il existe une courbe ferm\'ee int\`egre, localement principale $Y$ telle que $H={O}_X(Y)$. D'apr\`es le th\'eor\`eme de Riemann, (i') implique que $H^1\bigl(N_Y^{\otimes n}(p)\bigr)=0$ pour $p\geq0$ et $n\geq n_0$ avec $n_0=2h^1({O}_Y)-1$. Comme $H^2\bigl(N^{\otimes n}(p)\bigr)=0$ pour $p\gg0$, la suite exacte
\[
   H^1\bigl(N_Y^{\otimes n}(p)\bigr)
   \longrightarrow H^2\bigl(N^{\otimes n}(p-1)
   \longrightarrow H^2\bigl(N^{\otimes n}(p)\bigr)
   \longrightarrow0
\]
donne que $H^2(N^{\otimes n})=0$ pour $n\geq n_0$. Donc pour $n\gg0$, on a
\[
   h^0(N^{\otimes n})
   =\chi(N^{\otimes n})+h^1(N^{\otimes n})
   \geq\chi(N^{\otimes n})
   =1/2\langle c_1(N)^2\rangle n^2+\ldots>0,
\]
ce qui ach\`eve de prouver le lemme.

\textbf{7.1.3.} Maintenant supposons dans 7.1 qu'il existe un faisceau inversible $M$ sur $X$ tel que $\langle c_1(L)\cdot c_1(M)\rangle\ne0$, et raisonnons par l'absurde. Rempla\c{c}ons $X$ par le sch\'ema r\'eduit sous-jacent d'une pr\'esentation de Chow de $X$, ce qui est loisible gr\^ace \`a la formule de projection X 4.4.4, et soit $G$ ample sur $X$. Prenons $m>0$ tel que $H_1=G^{\otimes m}\otimes M$ soit ample, et tel que $\langle c_1(L)\cdot c_1(H_1)\rangle\ne0$ ; prenons $p,q$ ($q\ne0$) tels que, posant $L_1=L^{\otimes p}\otimes H^{\otimes q}$, on ait que $\langle c_1(L_1)\cdot c_1(H_1)\rangle=p\langle c_1(L)\cdot c_1(H_1)\rangle+q\langle c_1(H)\cdot c_1(H_1)\rangle$ est $0$ ; alors on a aussi que $\langle c_1(L_1)^2\rangle=p^2\langle c_1(L)^2\rangle+q^2\langle c_1(H)^2\rangle$ est $>0$ ; ce qui revient au cas 7.1.1, avec $L=L_1$ et $H=H_1$.

\textbf{7.1.4.} Enfin consid\'erons le cas g\'en\'eral de 7.1. Toujours raisonnant par l'absurde, supposons que $L$ ne soit pas num\'eriquement \'equivalent \`a $0$ ; en outre, rempla\c{c}ons $X$ par la normalis\'ee de $X_{\red}$, de sorte que $X$ n'a 
